\section{Síntesis y ampliación}

El núcleo de Switch Transformer no es «usar un router para obtener mágicamente un modelo de un billón de parámetros», sino plantear la expansión condicional de parámetros como un sistema de extremo a extremo: el router decide la ruta, la capacity restringe las rutas dinámicas a lotes estáticos, el load balance eleva la utilización del hardware, la selective precision y la inicialización protegen el entrenamiento, el expert dropout atiende el sobreajuste en tareas posteriores, la expert parallelism se encarga de distribuir los parámetros, y los experimentos deben separar step, time, FLOPs y costo total.

\begin{enumerate}
  \item \textbf{Lo que la dispersión escala es el conjunto de parámetros seleccionables}: los parámetros totales pueden crecer con el número de experts, mientras que los active experts por token siguen controlados por top-$k$;
  \item \textbf{Top-1 es una elección de Pareto}: reduce el FFN de expert, la comunicación y los requisitos de capacidad, pero que resulte superior depende del hardware y de la calidad objetivo;
  \item \textbf{Router es una ruta de control numérico}: softmax y argmax son sensibles al redondeo; usar float32 de forma local evita, con un costo mínimo, que la ruta se invierta;
  \item \textbf{Load balance es la interfaz entre el algoritmo y el sistema}: la loss auxiliar convierte probabilidades aprendibles en lotes por expert aproximadamente regulares;
  \item \textbf{Capacity factor es una perilla explícita de costo}: una capacidad mayor reduce el dropping, pero amplía el buffer, la comunicación y el cómputo potencial;
  \item \textbf{Los resultados negativos importan igual}: enviar los tokens de overflow a experts subóptimos no es necesariamente mejor que omitir la subcapa de expert;
  \item \textbf{Scaling debe validarse sobre varios ejes horizontales}: step, wall-clock, FLOPs, almacenamiento de parámetros y costo responden preguntas distintas;
  \item \textbf{Parallelism se puede combinar}: data, model y expert parallelism dividen, respectivamente, las muestras, los tensores y los parámetros condicionales;
  \item \textbf{La correlación en tareas posteriores no equivale a una prueba del mecanismo}: upstream loss, parámetros totales, active compute y cobertura de datos requieren un control emparejado;
  \item \textbf{La destilación es una conversión de despliegue con pérdida}: debe compararse con un student scratch de la misma escala y contabilizarse el costo de entrenamiento del teacher;
  \item \textbf{Vision priority routing avanza hacia el cómputo adaptativo}: cuando la capacity es menor que 1, el modelo decide activamente qué patches reciben cómputo de expert;
  \item \textbf{El escenario de aplicación determina el valor}: MoE se orienta más al entrenamiento a gran escala y a un rendimiento de procesamiento medio o alto; no es automáticamente adecuado para lotes pequeños, dispositivos de borde ni contextos extremadamente largos.
\end{enumerate}

\begin{importantbox}{Tarea del curso: realizar una auditoría matched-resource de MoE}
Implementa un Transformer dense pequeño y un MoE top-1, emparejando los tokens de entrenamiento, los active FLOPs y el wall-clock total. Barre al menos el número de experts, el capacity factor y el peso de la aux-loss; reporta PPL, time-to-quality, overflow, expert entropy, tiempo de all-to-all, memoria de GPU pico y per-task fine-tuning. Agrega además un dense baseline con el mismo número de parámetros y un dense baseline con los mismos active FLOPs, y responde si la ganancia proviene de la capacidad, del cómputo, de la regularización o de la implementación del sistema.
\end{importantbox}

\begin{warningbox}{Autoverificación final: seis conceptos que no deben confundirse}
Los \textbf{parámetros totales} no equivalen a los \textbf{parámetros activados}; los \textbf{FLOPs} no equivalen a la \textbf{latencia}; el \textbf{equilibrio de carga} no equivale a \textbf{expertos idénticos}; el \textbf{token dropping} no equivale a \textbf{eliminar tokens}; la \textbf{upstream correlation} no equivale al \textbf{efecto causal de los parámetros}; la \textbf{sparsity} no equivale a \textbf{cualquier adaptive computation}.
\end{warningbox}

\subsection{Lecturas adicionales}

Se sugiere leer los materiales originales en el orden «mecanismo histórico — expansión del sistema — receta de Switch — extensión multilingüe y visual»:

{\small
\begin{itemize}[nosep,leftmargin=*]
  \item \href{https://doi.org/10.1162/neco.1991.3.1.79}{\textbf{Adaptive Mixtures of Local Experts}}: primera formalización del gating y de los expertos locales.
  \item \href{https://arxiv.org/abs/1701.06538}{\textbf{Sparsely-Gated MoE}}: precursor clave del cómputo condicional a gran escala, el enrutamiento y el equilibrio de carga.
  \item \href{https://arxiv.org/abs/2001.08361}{\textbf{Scaling Laws for Neural Language Models}}: contexto empírico del dense scaling.
  \item \href{https://arxiv.org/abs/2006.16668}{\textbf{GShard}}: sharding automático y sistema de cómputo condicional a gran escala.
  \item \href{https://arxiv.org/abs/2101.03961}{\textbf{Switch Transformers}}: fuente principal sobre top-1 routing, selective precision, capacity, experimentos y destilación.
  \item \href{https://arxiv.org/abs/1910.10683}{\textbf{T5}} y \href{https://arxiv.org/abs/2010.11934}{\textbf{mT5}}: líneas base dense text-to-text y de 101 idiomas.
  \item \href{https://arxiv.org/abs/2106.05974}{\textbf{V-MoE}}: sparse experts visuales, priority routing y asignación adaptativa con capacidad menor que 1.
  \item \href{https://github.com/tensorflow/mesh}{\textbf{Mesh TensorFlow}}: contexto de la implementación distribuida y del código de sharding en que se desarrolló el trabajo original.
\end{itemize}
}

\end{document}
